\section{推荐阅读}

\begin{readingbox}
\textbf{基础}
\begin{itemize}[nosep]
  \item Eugene Izhikevich, \emph{Dynamical Systems in Neuroscience}，用于从生物学上具体的数学风格理解神经动力学、兴奋性类别、振荡与节律生成。
  \item Michael Brin and Garrett Stuck, \emph{Introduction to Dynamical
        Systems}，用于理解更一般数学层面上的不动点、稳定性、分岔与不变集思维。
  \item Steven Strogatz, \emph{Nonlinear Dynamics and Chaos}，用于在相图、极限环与阈值式定性变化之间建立最可读的桥梁。
  \item Arkady Pikovsky, Michael Rosenblum, and J\"urgen Kurths,
        \emph{Synchronization}，用于理解耦合、牵引与相位锁定。
\end{itemize}
\textbf{现代研究}
\begin{itemize}[nosep]
  \item 关于大脑中吸引子与积分器网络的 Nature Reviews Neuroscience
        (2022) 综述，作为本章主要的中等车道综述锚点。
  \item \emph{Metastability demystified: the foundational past, the pragmatic
        present and the promising future} （Nature Reviews Neuroscience, 2024,
        medium），用于理解灵活稳定性与重构。
  \item \emph{Falling asleep follows a predictable bifurcation dynamic}
        （Nature Neuroscience, 2025，在其自身生物学领域内属 strong），用于提供一个具体的转变模型，并为阈值敏感型推理提供正当性。
  \item \emph{Self-orthogonalizing attractor neural networks emerging from the
        free energy principle} （arXiv, 2025, medium），用于理解序列形成与自组织吸引子结构。
  \item \emph{Koopman-Nemytskii Operator: A Linear Representation of Nonlinear
        Controlled Systems} （arXiv, 2025, medium），用于作为算子与反馈视角的高级延伸。
\end{itemize}
\textbf{前沿}
\begin{itemize}[nosep]
  \item 高级算子方案与意识架构方案，应被视为形式化升级路径或假设生成器，而不是已完成的行为证明。
\end{itemize}
\end{readingbox}